\subsection{Memoria genealógica y contenido informacional de una lectura}
\label{sec:ii-memoria-lectura}

El operador de área reúne ocupaciones de frontera. El observable
ponderado añade su procedencia sectorial. Esta diferencia expresa un
principio físico del manuscrito: la información accesible depende de
la aplicación que extrae la coordenada observable del estado. En APP, el residuo y el cociente
conservan conjuntamente la operación entera; en TPK, la fase visible
convive con la orientación y el incremento de memoria. La realización
física transporta esa misma distinción entre estado, observable y
registro de recuperación \cite{informacion}.

Sea \(\Omega\) un conjunto finito de estados alcanzados por la
actualización considerada. Un lector \(q:\Omega\to Y\) asigna la
coordenada observable y \(M:\Omega\to\mathcal M\) conserva la memoria
de la fibra: hoja, orientación, cociente, acarreo y ruta pertinentes.
Para una distribución \(p\) sobre esos estados, denotemos por
\(\mathsf H_p\) la entropía adimensional, con logaritmo natural.

\begin{proposition}[Resolución informacional de la fibra]
Si \((q,M)\) es inyectiva sobre el soporte de \(p\), entonces
\begin{align}
 \mathsf H_p(X)&=\mathsf H_p(q(X))+
                 \mathsf H_p(X\mid q(X)),\\
 \mathsf H_p(X\mid q(X),M(X))&=0.
 \label{eq:ii-memoria-condicional}
\end{align}
\end{proposition}
\begin{proof}
La factorización de la distribución conjunta en distribución visible
y distribución condicional sobre la fibra da la primera igualdad.
La inyectividad proporciona una inversa sobre la imagen de \((q,M)\);
cada par expresado determina un único estado del soporte y su
distribución condicional está concentrada en ese estado.
\end{proof}

La ecuación \eqref{eq:recover} realiza este mecanismo para las dos
ocupaciones agregadas: \(N\) da el total y \(D\) resuelve su
distribución sectorial. Los estados \((N_{90},N_{120})=(1,0)\) y
\((0,1)\) comparten el área agregada y se distinguen por el segundo
observable. Bajo una distribución equiprobable sobre ese par, la
lectura exclusiva de \(N\) deja \(\log2\) nats de incertidumbre;
el par \((N,D)\) la resuelve exactamente. La afirmación concierne a
esas ocupaciones. Las microhistorias dentro de cada sector conservan
su registro adicional.

La interpretación de la memoria de borde queda así determinada por
una recuperación concreta. Una reducción puede ocultar información
que el estado conjunto sigue transportando. Para una actualización
biyectiva \(U\), la igualdad \(\mathsf H_p(X\mid U(X))=0\) expresa
reversibilidad lógica. Una operación física de descarte o reinicio
de memoria tiene, por su parte, su propio balance de información y
energía. El entrelazamiento se evalúa respecto de la bipartición y
el estado reducido especificados en la sección anterior.

\subsection{Hojas areal y volumétrica: incidencia, crecimiento y medida}
\label{sec:ii-hojas-av}

La distinción areal--volumétrica posee una realización dinámica en el
calendario del TPK. El bloque \((+1,-1,0)\), aplicado a los modos
\(\Sigma,\Pi\) mediante
\[
 +1:m\mapsto\Sigma,\qquad -1:m\mapsto\Pi,\qquad 0:m\mapsto m,
\]
produce la órbita cíclica \((\Sigma,\Pi,\Pi)\). El lector geométrico
realiza las correspondencias \(\Sigma\mapsto\mathscr H_{\rm ar}\) y
\(\Pi\mapsto\mathscr H_{\rm vol}\). Sus transiciones consecutivas
fijan la incidencia
\begin{equation}
 F_{\rm av}=\begin{pmatrix}0&1\\1&1\end{pmatrix},\qquad
 N_9=3F_{\rm av},\quad N_{27}=9F_{\rm av},\quad
 N_{108}=36F_{\rm av}.
 \label{eq:ii-hojas-incidencia}
\end{equation}
En efecto, el ciclo contiene una transición de cada tipo
\(\Sigma\to\Pi\), \(\Pi\to\Pi\), \(\Pi\to\Sigma\).
Los calendarios mayores repiten ese bloque con las multiplicidades
indicadas.

El envolvente simbólico asociado permite todas las concatenaciones
admitidas por \(F_{\rm av}\). Su número de caminos depende de las
potencias de esa matriz. Con \(F_0=0\), \(F_1=1\), la inducción da
\begin{equation}
 F_{\rm av}^{\,n}=
 \begin{pmatrix}F_{n-1}&F_n\\F_n&F_{n+1}\end{pmatrix},
 \qquad h_{\rm top}=\log\varphi.
 \label{eq:ii-hojas-fibonacci}
\end{equation}
La identidad matricial usa la recurrencia de Fibonacci. El crecimiento
exponencial queda fijado por la raíz de Perron \(\varphi\) de la
incidencia ya construida.

\begin{proposition}[Distribución estacionaria de las hojas]
La medida de máxima entropía del envolvente simbólico tiene pesos
\begin{equation}
 \mu_{\rm ar}=\frac{1}{1+\varphi^2},\qquad
 \mu_{\rm vol}=\frac{\varphi^2}{1+\varphi^2},\qquad
 \frac{\mu_{\rm ar}}{\mu_{\rm vol}}=\varphi^{-2}.
 \label{eq:ii-hojas-medida}
\end{equation}
\end{proposition}
\begin{proof}
El vector positivo \(r=(1,\varphi)\) satisface
\(F_{\rm av}r=\varphi r\). Las probabilidades de transición
\(P_{ij}=(F_{\rm av})_{ij}r_j/(\varphi r_i)\) suman uno.
La simetría de \(F_{\rm av}\) da pesos estacionarios proporcionales
a \(r_i^2\). Su normalización produce
\eqref{eq:ii-hojas-medida}. En la media estacionaria,
\(-\log P_{ij}=\log\varphi+\log r_i-\log r_j\) sobre las
aristas permitidas. Los dos últimos términos se cancelan y la
entropía por transición es \(\log\varphi\), igual a la entropía
topológica de \eqref{eq:ii-hojas-fibonacci}.
\end{proof}

El calendario periódico y su envolvente simbólico constituyen dos
objetos explícitos: el primero posee frecuencias cronológicas
\((1/3,2/3)\); el segundo admite las concatenaciones de la incidencia
y posee la medida \eqref{eq:ii-hojas-medida}. Esta separación permite
atribuir a cada cantidad su significado físico-matemático. La razón
\(\varphi^{-2}\) caracteriza la distribución de las hojas en esa
medida; \(\log\varphi\) caracteriza el crecimiento de sus historias
permitidas. La lectura angular de frontera incorpora después la
coordenada regional correspondiente. La construcción distingue
crecimiento volumétrico, representación areal y memoria del transporte.

\subsection{Continuaciones admisibles e información de supervivencia}
\label{sec:ii-informacion-supervivencia}

Fijado un estado TPK \(x\), sea \(\Omega_n(x)\) el conjunto finito
de sus continuaciones de longitud \(n\) admitidas por la regla de
prolongación. Las restricciones entre niveles conservan los prefijos.
Una colección compatible de probabilidades \(p_n\), obtenida del
lector estadístico que se declare, define la incertidumbre de
continuación
\begin{equation}
 s_n(x)=-\sum_{\omega\in\Omega_n(x)}
                 p_n(\omega)\log p_n(\omega).
 \label{eq:ii-entropia-continuaciones}
\end{equation}
La compatibilidad exige que la probabilidad de un prefijo sea la suma
de las probabilidades de sus prolongaciones. En el envolvente anterior,
la medida estacionaria proporciona una colección de este tipo.

\begin{proposition}[Balance informacional del refinamiento]
Para la colección compatible anterior,
\begin{align}
 0\leq s_n(x)&\leq\log|\operatorname{supp}p_n|,\\
 s_{n+1}(x)-s_n(x)
 &=\mathsf H(X_{n+1}\mid X_n)\geq0.
 \label{eq:ii-entropia-refinamiento}
\end{align}
La cota superior se alcanza para la distribución uniforme sobre el
soporte. El incremento es nulo cuando cada prefijo de probabilidad
positiva determina su prolongación con probabilidad uno.
\end{proposition}
\begin{proof}
La desigualdad elemental para la entropía de una distribución finita
da la primera línea. Como \(X_n\) es la restricción de \(X_{n+1}\),
la factorización de probabilidades sobre cada fibra de prolongación
da \(\mathsf H(X_{n+1})=\mathsf H(X_n)+
\mathsf H(X_{n+1}\mid X_n)\). Cada entropía condicional es no
negativa y se anula exactamente para una distribución puntual.
\end{proof}

Las continuaciones que permanecen admisibles tienen así un contenido
informacional preciso. Su multiplicidad y sus pesos describen la
incertidumbre residual de la lectura a profundidad finita. Una
historia realizada conserva, además, su orientación y memoria. La
entropía informacional, la entropía del estado reducido y la entropía
termodinámica se enlazan mediante sus respectivos lectores; cada una
conserva la información de dominio necesaria para su interpretación.
